\section*{Hoofdstuk \ref{ch:g8:hormonal-communication} --- Communicatie via hormonen}

\begin{solution}{exo:g8:hormonal-communication:1}
Dat ze haar product in het \omterm{prop:g4:journey-of-food:absorb}{bloed} afgeeft en niet door een buisje. De
regelklier onder in de \omterm{def:g5:brain-in-command:system}{hersenen}: dirigent van de andere en starter van de
\omterm{def:g8:puberty:puberty}{puberteit}. De schildklier: de tempoklier van het lichaam. De bijnieren:
het alarmhormoon. De alvleesklier: het beheer van de bloedglucose.
(\omterm{def:g8:reproductive-systems:female}{Eierstokken} en \omterm{def:g8:reproductive-systems:male}{zaadballen}: de cyclus- en verbouwingshormonen.)
\end{solution}

\begin{solution}{exo:g8:hormonal-communication:2}
Het vervoer is een omroep, de adressering zit bij de ontvanger: alleen
cellen met de receptoren van dat \omterm{def:g8:hormonal-communication:glands}{hormoon} reageren. Elk huis krijgt de
brief; alleen de passende sleutel opent hem.
\end{solution}

\begin{solution}{exo:g8:hormonal-communication:3}
Het tempo van het \omterm{def:g4:heart-and-blood:heart}{hart} omhoog (\cref{ch:g7:heart-and-circulation}); de
luchtwegen wijder (\cref{ch:g7:lungs-and-gas-exchange}); de glucosebank van
de \omterm{prop:g7:digestion-and-nutrients:absorption}{lever} vrijgegeven
(\cref{ch:g7:digestion-and-nutrients}); het \omterm{prop:g4:journey-of-food:absorb}{bloed} omgeleid naar de \omterm{def:g3:bones-and-muscles:muscle}{spieren}
(\cref{ch:g7:organs-at-work}) --- met de bleke, prikkende \omterm{def:g1:the-five-senses:sense}{huid} die de reeks
compleet maakt.
\end{solution}

\begin{solution}{exo:g8:hormonal-communication:4}
De schrijvende \omterm{def:g8:hormonal-communication:glands}{klier} houdt het resultaat van haar brieven in de gaten en
vertraagt naarmate het doel bereikt wordt --- thermostaat, geen kraan.
Paradepaardje: de bloedglucose, bij het stijgen gebankierd en bij het dalen
vrijgegeven door het hormoonpaar van de alvleesklier.
\end{solution}

\begin{solution}{exo:g8:hormonal-communication:5}
Het signaal van het opslaghormoon --- niet geschreven, of niet beantwoord
door zijn ontvangers. De \omterm{prop:g7:organs-at-work:bill}{glucose} stapelt zich daarna na de maaltijden op,
overweldigt het terugwinnen door de \omterm{def:g7:removing-wastes:kidneys}{nier}, en verschijnt in de \omterm{def:g7:removing-wastes:kidneys}{urine}.
\end{solution}

\begin{solution}{exo:g8:hormonal-communication:6}
Onder in de \omterm{def:g5:brain-in-command:system}{hersenen} --- de regelklier, met het \omterm{def:g5:brain-in-command:system}{zenuwstelsel} erboven en de
\omterm{def:g8:hormonal-communication:glands}{hormonen} eronder. Voorbeeld: de schrik --- eerst zenuwsignalen, daarna de
brief van de bijnier die het alarm volhoudt; of de daglengte die in de
voortplantingshormonen van de schapen wordt omgezet.
\end{solution}

\begin{solution}{exo:g8:hormonal-communication:7}
Een kraan giet tot iemand haar dichtdraait; een thermostaat meet wat hij
verandert en stelt zichzelf bij. De \omterm{def:g8:hormonal-communication:glands}{klier} leest het gehalte af dat ze
regelt --- de afgifte stijgt en zakt met de behoefte mee, zonder dat er een
hand van buitenaf nodig is.
\end{solution}

\begin{solution}{exo:g8:hormonal-communication:8}
Schrijver: de alvleesklier, aangezet door stijgende bloedglucose.
Ontvangers: vooral de \omterm{prop:g7:digestion-and-nutrients:absorption}{lever} en de \omterm{def:g3:bones-and-muscles:muscle}{spieren} --- die bankieren het overschot
bij bezorging. \omterm{prop:g8:hormonal-communication:feedback}{Terugkoppeling}: het dalende gehalte laat de afgifte zelf
wegzakken. Falen en gebruik: \omterm{ex:g8:hormonal-communication:diabetes}{diabetes} --- en de afgemeten doses van de
geneeskunde, de brief per injectiespuit.
\end{solution}

\begin{solution}{exo:g8:hormonal-communication:9}
Allebei schrijven ze met de hand een brief van het lichaam: de spuit
levert de ontbrekende opslagboodschap van de alvleesklier bij de
maaltijden; de pil levert dagelijks de boodschap ``opgeschort'' van de
cyclus. Dezelfde handeling --- een \omterm{def:g8:hormonal-communication:glands}{hormoon} nabootsen om zijn stelsel met
opzet te draaien --- de een herstelt een uitgevallen briefwisseling, de
ander overstemt met opzet een gezonde.
\end{solution}

\begin{solution}{exo:g8:hormonal-communication:10}
Het begin van het rillen: \omterm{def:g5:brain-in-command:system}{zenuwen} --- spierbevelen in een fractie van een
seconde. De baard: \omterm{def:g8:hormonal-communication:glands}{hormonen} --- maanden, het hele lichaam, de bloedbaan.
De glucosebaan: \omterm{def:g8:hormonal-communication:glands}{hormonen} --- een briefwisseling met \omterm{prop:g8:hormonal-communication:feedback}{terugkoppeling}. De
kniepeesreflex: \omterm{def:g5:brain-in-command:system}{zenuwen} --- één koppeling, een vijftigste seconde.
\end{solution}

\begin{solution}{exo:g8:hormonal-communication:11}
Het waarnemen van korter wordende dagen door de \omterm{def:g1:the-five-senses:sense}{ogen} gaat via de \omterm{def:g5:brain-in-command:system}{zenuwen};
het kantoor onder in de \omterm{def:g5:brain-in-command:system}{hersenen} vertaalt dat en dirigeert de regelklier;
haar brieven wekken het eigen hormoonprogramma van de \omterm{def:g8:reproductive-systems:female}{eierstokken}; en de
cycli van de ooien beginnen --- licht erin bij de telegraaf, lammeren eruit
per post.
\end{solution}

\begin{solution}{exo:g8:hormonal-communication:12}
Verdediging: één \omterm{prop:g4:journey-of-food:absorb}{bloed} draagt elke brief, en toch werkt elke brief alleen
waar receptoren antwoorden --- het publiek bepaalt het kanaal, dus wordt
een omroep een net van privélijnen. De doorzetting: een \omterm{def:g5:first-look-at-cells:cell}{cel} voegt zich bij
een publiek door de receptor te bouwen --- dus welke weefsels luisteren
wordt zelf gebouwd, orgaan voor orgaan, tijdens de ontwikkeling; de
doelwitten van de \omterm{def:g8:puberty:puberty}{puberteit} --- stembandenkast, \omterm{def:g1:the-five-senses:sense}{huid}, groeischijven ---
kregen hun ontvangers vooraf ingebouwd en wachtten jaren op de eerste brief.
\end{solution}

\begin{solution}{pb:g8:hormonal-communication:1}
\textbf{1.} Het briefpaar van de alvleesklier --- het opslaghormoon bij het
stijgen, zijn vrijgevende partner bij het dalen --- dat de baan met
\omterm{prop:g8:hormonal-communication:feedback}{terugkoppeling} vasthield terwijl de eigenaar sliep.

\textbf{2.} Sein: de \omterm{prop:g7:organs-at-work:bill}{glucose} van de maaltijd die de \omterm{prop:g7:digestion-and-nutrients:absorption}{darmvlokken} oversteekt
en het bloedgehalte optilt. Schrijver: de alvleesklier. Brief: het
opslaghormoon. Ontvangers: de \omterm{prop:g7:digestion-and-nutrients:absorption}{lever} en de \omterm{def:g3:bones-and-muscles:muscle}{spieren}, die het overschot
bankieren.

\textbf{3.} Een kraan zou voorbij het doel doorgieten en tot een duik
leiden; dat de daling in de baan afvlakt, laat zien dat de afgifte wegzakt
naarmate haar eigen resultaat binnenkomt --- zelfmetend, zelfstoppend.

\textbf{4.} Dat ze met het gehalte meesteeg en er weer mee wegstierf: een
curve van afgifte die de curve van de \omterm{prop:g7:organs-at-work:bill}{glucose} schaduwt, met de piek net
achter de piek --- het verloop van de thermostaat.

\textbf{5.} De vrijgevende partner (met hulp van het alarmhormoon bij
zware inspanning): de gebankierde \omterm{prop:g7:organs-at-work:bill}{glucose} van de \omterm{prop:g7:digestion-and-nutrients:absorption}{lever} --- de bank die bij
het ontbijt gevuld werd op de manier van
\cref{ex:g7:digestion-and-nutrients:breadtrace} --- uitgegeven aan de
rekening van de werkende \omterm{def:g3:bones-and-muscles:muscle}{spieren}.

\textbf{6.} Diezelfde vrijgevende partner: die de bank gestaag door de
voedselloze uren heen ontsluit en het uitgeven van het lichaam bijhoudt ---
de baan van de andere kant vastgehouden.

\textbf{7.} De bijnieren. De suiker legt brandstof klaar voor
\emph{vluchten of vechten}: het hele punt van de alarmbrief is een lichaam
dat bevoorraad en omgeleid is voor onmiddellijke inspanning --- of die
inspanning nu komt of niet.

\textbf{8.} Omdat de gezonde briefwisseling erop antwoordde: het
vrijgegeven overschot, dat niet uitgegeven werd (de bijna-aanrijding ging
voorbij), zette op zijn beurt de opslagbrief aan, en de baan werd hersteld
--- alarm en thermostaat achter elkaar.

\textbf{9.} Onbehandeld: de stijging na het ontbijt klimt door en daalt
niet --- er wordt geen opslagbrief geschreven of gehoord --- een hoog
plateau de hele ochtend; en voorbij de terugwindrempel verschijnt er suiker
in de \omterm{def:g7:removing-wastes:kidneys}{urine} (het signaal uit \cref{ex:g7:removing-wastes:thrift}).

\textbf{10.} De eigen maaltijdbrief van de alvleesklier. De dosis moet bij
het bord passen omdat de grootte van de brief moet passen bij de \omterm{prop:g7:organs-at-work:bill}{glucose}
die ze moet bankieren --- te klein laat een opeenhoping staan, te groot
schiet de baan naar beneden voorbij, en dat is op zichzelf een noodgeval.

\textbf{11.} 's Nachts en 's middags: het paar van de alvleesklier, met
\omterm{prop:g7:digestion-and-nutrients:absorption}{lever} en \omterm{def:g3:bones-and-muscles:muscle}{spieren} die antwoorden --- de baan vastgehouden. Bij het ontbijt
en het avondeten: de \omterm{prop:g7:organs-at-work:bill}{glucose} spreekt, de alvleesklier schrijft, de \omterm{prop:g7:digestion-and-nutrients:absorption}{lever}
bankiert. 's Middags: de inspanning geeft uit, de vrijgevende partner
ontsluit. Om vijf uur: de bijnieren onderbreken, de alvleesklier ruimt
daarna op. Eén waarde, vijf sprekers, de hele dag geen stilte.

\textbf{12.} ``Een stelsel met \omterm{prop:g8:hormonal-communication:feedback}{terugkoppeling}: schrijvers die de
resultaten van hun eigen brieven lezen --- en zo rijdt een waarde een leven
lang in een smalle baan zonder dat er iemand stuurt.''
\end{solution}
